\documentclass[11pt,a4paper]{article}
\usepackage[margin=25mm]{geometry}
\usepackage{amsmath,amssymb,amsthm,booktabs,hyperref}
\newtheorem{theorem}{Theorem}
\newtheorem{lemma}[theorem]{Lemma}
\newtheorem{proposition}[theorem]{Proposition}
\title{Quaternionic characters in a vertex permutation representation}
\author{Human Compression project: recovery of L024}
\date{}
\begin{document}\maketitle
We recover the example and infinite-family conclusion of Cameron,
Herman and Leemans~\cite{CHL}. The example answers Monson's existence
question and its polytopal variant; these answers belong to that source.
Our account fixes explicit matrices, identifies the subgroup fusion
without a character-table search, and uses a common character bound
to shorten the positivity verification. It also separates an erroneous
small-dimensional case in the printed infinite-family statement.

\section{Fixed vectors and a dihedral stabilizer}
Let $G$ be finite and $H\le G$. For an irreducible complex character
$\chi$, Frobenius reciprocity and the averaging projection give
\begin{equation}\label{average}
 [1_H^G,\chi]=\dim V_\chi^H
 =\frac1{|H|}\sum_{h\in H}\chi(h).
\end{equation}
Indeed $|H|^{-1}\sum_{h\in H}\rho(h)$ is an idempotent with image
the fixed subspace, and its trace is the right side. The permutation
module on $G/H$ has that same multiplicity by reciprocity. For a
vertex stabilizer in a regular polytope, the nontrivial constituents
are precisely the polytopal characters in the terminology of
\cite{CHL}. The trivial constituent is removed in the deleted
permutation module.

\begin{lemma}\label{dihedral}
Suppose $H$ is dihedral of order $2\ell$, where $\ell$ is an odd prime.
Assume all its reflections lie in a $G$-class $t^G$ and all its
nonidentity rotations lie in a $G$-class $u^G$. Then
\[
 [1_H^G,\chi]=\frac{\chi(1)+\ell\chi(t)+(\ell-1)\chi(u)}{2\ell}.
\]
In particular this multiplicity is positive whenever
\[
 \chi(1)>\ell\sqrt{|C_G(t)|}+(\ell-1)\sqrt{|C_G(u)|}.
\]
\end{lemma}
\begin{proof}
There are $\ell$ reflections and $\ell-1$ nonidentity rotations,
so the formula follows from~\eqref{average}. Column orthogonality
gives $\sum_{\psi\in\operatorname{Irr}(G)}|\psi(x)|^2=|C_G(x)|$.
Consequently $|\chi(x)|\le\sqrt{|C_G(x)|}$. Applying the triangle
inequality to the numerator yields the sufficient condition. The
average is a nonnegative integer by~\eqref{average}.
\end{proof}

We also retain the Frobenius--Schur theorem: the indicator
\[
 \nu_2(\chi)=|G|^{-1}\sum_{x\in G}\chi(x^2)
\]
is $-1$ precisely for a real-valued irreducible character of real
Schur index two. These are standard character-theoretic inputs,
not new theorems of this work.

\section{A fixed marking of \texorpdfstring{$\operatorname{PSU}(5,3)$}{PSU(5,3)}}
Let $E=\mathbb F_3[z]/(z^2-z-1)$, let $J$ be the $5$ by $5$
antidiagonal identity, and write $\bar x=x^3$. We use row vectors
and the unitary condition $sJ\bar s^{,t}=J$. Define
\[
 s_0=\operatorname{diag}(-1,1,1,1,-1),\qquad
 s_2=\begin{pmatrix}
 0&0&0&0&1\\0&1&0&0&0\\0&0&-1&0&0\\0&0&0&1&0\\1&0&0&0&0
 \end{pmatrix},
\]
and
\[
 s_1=\begin{pmatrix}
 z^7&z^3&z^7&0&z^4\\
 z^5&z^6&0&z^4&0\\
 z^3&z^4&z^4&0&z^5\\
 z^5&z^4&z^4&z^2&z\\
 0&z^7&z&z^7&z^5
 \end{pmatrix}.
\]
The center of $\operatorname{SU}(5,3)$ has order $\gcd(5,4)=1$,
so this group equals $G=\operatorname{PSU}(5,3)$. Its standard
unitary order formula gives $|G|=258190571520$.

\begin{proposition}\label{marking}
The ordered triple $(s_0,s_1,s_2)$ is a string C-group marking of
$G$ of type $\{5,5\}$. Its vertex stabilizer $H=\langle s_1,s_2\rangle$
is dihedral of order ten.
\end{proposition}
\begin{proof}
Direct arithmetic in $E$ verifies that all three matrices have
determinant one, preserve $J$, and are involutions. It gives
\[
 |s_0s_1|=|s_1s_2|=5,\qquad |s_0s_2|=2.
\]
For the finite generation step we use the exact attached certificate,
\texttt{verify\_schur\_example.g}. It enumerates the projective
isotropic points by normalizing the first nonzero coordinate to one
and testing $\sum_i v_i v_{6-i}^3=0$. There are $2440$ points.
The three displayed matrices act on this literal set. A stabilizer
chain for their permutation group has successive orbit lengths
\[
 2440,\quad2187,\quad2,\quad224,\quad36,\quad3.
\]
Their product is $258190571520$. Thus the subgroup generated by the
matrices has an image of order $|G|$, and containment in $G$ proves
equality. The verification needs no search for the generator triple or
downloaded group representation. The permutation stabilizer-chain
option \texttt{random=1000} requires guaranteed correctness of the
verification, even though the internal construction may use random
choices; the certificate asserts this setting. The chain itself is recomputed exactly from the fixed
matrices; the displayed product alone would not establish generation.

We now prove the intersection property from the product orders and
full generation. Two involutions whose product has order five generate
a dihedral group of order ten. Thus the adjacent parabolics
$P_{01}=\langle s_0,s_1\rangle$ and
$P_{12}=\langle s_1,s_2\rangle$ are both dihedral of order ten,
whereas $P_{02}=\langle s_0,s_2\rangle$ is a Klein four-group.
A subgroup of a dihedral group of order ten that contains a reflection
has order two or ten. Since $P_{01}\cap P_{12}$ contains $s_1$,
a larger intersection would make the two parabolics equal and put
all three generators in a group of order ten. Full generation excludes
this, so their intersection is $\langle s_1\rangle$.
The intersection of $P_{02}$ with either adjacent parabolic has order
dividing $\gcd(4,10)=2$; its common displayed generator therefore
generates the whole intersection. A singleton generator cannot belong
to its opposite parabolic, since that too would put all generators in
a proper parabolic. Distinct singleton subgroups meet trivially, and
all containment, empty-set and whole-group cases are immediate.
Consequently for every $I,K\subseteq\{0,1,2\}$,
\[
 \langle s_i:i\in I\rangle\cap\langle s_i:i\in K\rangle
 =\langle s_i:i\in I\cap K\rangle.
\]
The certificate also checks all $64$ pairs as corroboration; that part
of the computation is not a premise of this intersection argument.
The generation computation remains a finite premise, not a Lean
formalization.
\end{proof}

The source's two centralizer choices can also be identified directly.
An involution of $G$ has a nondegenerate negative eigenspace of
dimension two or four. The corresponding class sizes are $444690$
and $4941$, with centralizer orders $580608$ and $52254720$.
Our three matrices have negative eigenspace dimension two.
For $s_0$, write its orthogonal decomposition as $A\perp B$, with
$\dim A=2$, $\dim B=3$. Then
\[
 C_G(s_0)=\{(a,b)\in\operatorname{GU}(A)\times
 \operatorname{GU}(B):\det(a)\det(b)=1\}.
\]
The negative eigenspace dimensions $(k,l)$ of an involution in this
centralizer must have positive even sum. The five possibilities are
$(0,2),(1,1),(1,3),(2,0),(2,2)$. Their class sizes are
$63,378,6,1,63$. To see this, unitary groups are transitive on
nondegenerate subspaces of each dimension, and the numbers of
nondegenerate lines in dimensions two and three over $E$ are
\[
 \frac{|\operatorname{GU}(2,3)|}{|\operatorname{GU}(1,3)|^2}=6,
 \qquad
 \frac{|\operatorname{GU}(3,3)|}
 {|\operatorname{GU}(1,3)|\,|\operatorname{GU}(2,3)|}=63.
\]
The determinant-one constraint does not split these orbits: the
centralizer of either eigenspace decomposition has a determinant
surjective unitary factor, which corrects any conjugator's determinant.
The matrix $s_2$ has $(k,l)=(1,1)$, so it lies in the prescribed
largest class of size $378$.

\section{The character calculation}
We retain the ordinary character table of $G$ distributed by the
GAP Character Table Library. It has $150$ irreducibles, the two
involution classes above, and exactly one class of elements of order
five, whose centralizer has order $80$. Call the larger involution
class $2B$ and the order-five class $5A$. All five reflections of
$H$ are conjugate within $H$ and hence lie in $2B$, while its four
nonidentity rotations lie in $5A$. This determines the fusion needed
for induction without choosing one among possible table fusions.
Lemma~\ref{dihedral} gives
\begin{equation}\label{multiplicity}
 m_\chi=\frac{\chi(1)+5\chi(2B)+4\chi(5A)}{10}.
\end{equation}

Since $\sqrt{580608}<762$ and $\sqrt{80}<9$, every irreducible
of degree greater than $3846$ has positive multiplicity. The retained
table has precisely $24$ characters below this threshold, and the
smallest remaining degree is $4270$. For those $24$ characters all
needed data are displayed below; equal rows are grouped in the first
column.
\begin{center}
\begin{tabular}{rrrrr}\toprule
Number&$\chi(1)$&$\chi(2B)$&$\chi(5A)$&$m_\chi$\\\midrule
1&1&1&1&1\\1&60&4&0&8\\1&61&5&1&9\\2&61&9&1&11\\
1&549&-11&-1&49\\1&610&-6&0&58\\2&610&22&0&72\\
1&1220&-12&0&116\\2&1220&12&0&128\\
1&1281&49&1&153\\2&1281&21&1&139\\
1&1830&38&0&202\\2&1830&10&0&188\\
1&1890&42&0&210\\2&2440&24&0&256\\
1&3660&20&0&376\\2&3660&-36&0&348\\\bottomrule
\end{tabular}
\end{center}
Equation~\eqref{multiplicity} gives the last column. The remaining
multiplicities exceed $(4270-3846)/10>42$. Hence every nontrivial
irreducible is polytopal, and the least such multiplicity is eight,
attained at degree sixty.

The certificate computes the indicators directly from the class sizes
and the square power map, using their defining sum. Exactly five are
$-1$. Their degrees and multiplicities are
\[
\begin{array}{c|rrrrr}
\chi(1)&60&3660&7320&10980&14580\\\hline
m_\chi&8&376&752&1072&1440.
\end{array}
\]
The character table and its power map are borrowed data. Their
construction is not reproved here. The fixed-matrix generation certificate and written intersection proof,
together with these exact table computations,
recovers the source's affirmative answers to both existence questions.

\section{The infinite family and its small dimensions}
We retain the two character-theoretic results quoted in the source:
Lusztig's criterion for the existence of the cuspidal unipotent
character at triangular dimension $n=s(s+1)/2$, and Ohmori's
criterion that its real Schur index is two precisely when
$\lfloor n/2\rfloor$ is odd, in odd characteristic. We also retain
the quoted Nuzhin generation theorem in its stated range $n\ge8$.

For completeness about the two smaller triangular dimensions, the
subsequent primary classification~\cite[Theorem 1.1 and following
paragraph]{LV} says that $\operatorname{PSU}(3,q)$ admits no string
C-group representation, whereas $\operatorname{PSU}(6,q)$ admits
one of rank three for odd $q$. These are expressly borrowed
classification inputs here. This use does not constitute a recovery
of that separate classification paper.

\begin{theorem}
Let $q$ be an odd prime power, and let $n=s(s+1)/2$ with
$\lfloor n/2\rfloor$ odd. Then $\operatorname{PSU}(n,q)$ has a
string C-group representation and the indicated irreducible character
of real Schur index two if $n\ge6$. When $n=3$ it has no string
C-group representation at any rank.
\end{theorem}
\begin{proof}
The triangular dimensions below eight are $1,3,6$. The parity
condition excludes $1$. The borrowed classification settles $3$
and $6$, and the stated Nuzhin theorem supplies rank three for all
remaining admissible dimensions. Lusztig and Ohmori give the
character in every positive case. There are infinitely many such
dimensions: put $s=8j+5$, so $n=32j^2+44j+15\equiv3\pmod4$.
\end{proof}

The printed Theorem 4 of~\cite{CHL} omits the dimensional
qualification, and therefore includes the false $n=3$ case. The
infinite-family conclusion itself survives; the $n=6$ extension uses
the additional input explicitly identified above. The source's final
suggestion that these cuspidal characters should always be polytopal
is not established by the present argument and remains a separate
question in this recovery.

\section{Scope of the recovery}
The source reports that a particular implementation, choosing one
representative from each of $150$ conjugacy classes and conjugating
$s_0$ by it, returned $36$ markings, one of type $\{5,5\}$. Those
representatives and that run are not specified in the manuscript;
the number is an implementation report, not an invariant enumeration
of all markings. We neither reproduce that historical run nor promote
it to such a classification. Proposition~\ref{marking} supplies the
fully specified witness needed for the principal existence and
character conclusions. Its matrices replace the source's unspecified
choices. The general fixed-vector formula, column bound, source
attributions, corrected parameter range and computational premises
are all explicit. No new solution of an open conjecture, formal Lean
proof, or compression measurement is claimed by this manuscript.

\begin{thebibliography}{9}
\bibitem{CHL} P. J. Cameron, A. Herman and D. Leemans,
\emph{String C-groups with real Schur index 2},
arXiv:2201.01313v1 (4 January 2022), six pages;
J. Pure Appl. Algebra \textbf{226} (2022), 107025.
The inspected source is the stated preprint.
\bibitem{LV} D. Leemans and A. Vandenschrick,
\emph{On simple groups that admit a string C-group representation},
retained six-page author manuscript, Theorem 1.1 and its following
rank-three paragraph. The general classification is used as an input.
\bibitem{CT} The GAP Group and T. Breuer,
\emph{GAP Character Table Library}, version 1.3.11,
ordinary table \texttt{U5(3)}; GAP 4.15.1.
\end{thebibliography}
\end{document}
